\section{Introduction}

The fundamental concept for this study, and for interpreting the type of data produced by the reaction code TALYS, is the \emph{nuclear cross section}. Cross sections quantify reaction probabilities. The energy spectrum of the emitted particles corresponds to the differential cross section with respect to energy. For emission of a particle \(b\) into a solid-angle element \(d\Omega\), the differential cross section \(d\sigma/d\Omega\) has units of barn/sr. When the outgoing energy of the emitted particle, denoted here by \(E_{\mathrm{out}}\), is resolved, the relevant observable is the double-differential quantity
%
\begin{equation}
\frac{d^2\sigma}{d\Omega, dE_{\mathrm{out}}}
\label{eq:ddxs}
\end{equation}
%
which describes how the probability of emission is distributed in angle and outgoing energy.

In this project, nuclear reactions are modelled with TALYS and the resulting energy-differential light-ion spectra are compared with selected experimental datasets from EXFOR \cite{koning2025talys,iaea_exfor}. TALYS can simulate nuclear reactions by combining several physical mechanisms within a consistent reaction-chain description \cite{koning2025talys}. An incident projectile on a target nucleus may lead to a sequence of residual nuclei, each requiring appropriate nuclear-structure and reaction inputs, such as separation energies, optical potentials, and level densities. For each reaction, TALYS provides both the total predicted spectrum and a decomposition into contributions associated with different mechanisms: direct-like, pre-equilibrium, compound, and multiple pre-equilibrium emission \cite{koning2025talys}.
